% Fig. 1.12: degree-six Bezier evaluated in Bernstein form, not a cubic fit.
% P=[(-2,0),(-3,2),(1,3),(3,1),(2,-2),(-1,-2),(-2,0)].
% P0=P6 and P1-P0=P6-P5=(-1,2): position AND tangent match.
% Numeric control data are representative; the original gives no coordinates.
\begin{tikzpicture}[book line,line cap=round,line join=round,scale=\BookGeometryScale]
  \foreach \i/\x/\y in {0/-2/0,1/-3/2,2/1/3,3/3/1,4/2/-2,5/-1/-2,6/-2/0}
    \coordinate (P\i) at (\x,\y);
  \foreach \i/\j in {0/1,1/2,2/3,3/4,4/5,5/6}
    \draw[line width=.45pt,shorten <=2pt,shorten >=2pt] (P\i)--(P\j);
  \draw[shorten <=2pt,shorten >=2pt] plot[variable=\u,domain=0:1,samples=161]
    ({-2*(1-\u)^6-18*\u*(1-\u)^5+15*\u^2*(1-\u)^4
        +60*\u^3*(1-\u)^3+30*\u^4*(1-\u)^2-6*\u^5*(1-\u)-2*\u^6},
     {12*\u*(1-\u)^5+45*\u^2*(1-\u)^4+20*\u^3*(1-\u)^3
        -30*\u^4*(1-\u)^2-12*\u^5*(1-\u)});
  \node[book point,label=left:{$\mathbf P_0=\mathbf P_6$}] at (P0) {};
  \foreach \i/\pos in {1/above left,2/above,3/right,4/below,5/below}
    \node[book point,label=\pos:{$\mathbf P_{\i}$}] at (P\i) {};
\end{tikzpicture}
